% ===================================================================
% Table: Ablation study summary (3 dimensions, real training data)
% ===================================================================
\begin{table}[htbp]
  \centering
  \caption{消融实验汇总}
  \label{tab:ablation_summary}
  \small
  \setlength{\tabcolsep}{12pt}
  \begin{tabular}{lcc}
  \hline
  消融维度与变体 & 宏平均 F1 & 相对完整模型 \\
  \hline
  \textbf{自注意力模块} & & \\
  \quad 无 (CNN-only) & 0.884 & $-$0.5\% \\
  \quad 有 (CNN+Attention, \textbf{本文}) & \textbf{0.889} & — \\
  \hline
  \textbf{窗口大小 $T$} & & \\
  \quad $T=16$ & 0.840 & $-$5.5\% \\
  \quad $T=32$ (\textbf{本文}) & \textbf{0.889} & — \\
  \quad $T=48$ & 0.876 & $-$1.5\% \\
  \hline
  \textbf{训练数据规模} & & \\
  \quad 10\%（$\approx$1名参与者） & 0.790 & $-$11.1\% \\
  \quad 50\%（$\approx$5名参与者） & 0.865 & $-$2.7\% \\
  \quad 100\%（\textbf{本文}） & \textbf{0.889} & — \\
  \hline
  \end{tabular}
  \noindent\begin{minipage}{\linewidth}\footnotesize
  注：固定其他参数不变，仅改变所标注的消融维度。自注意力贡献相对有限（+0.5\%），根因在于手背类瓶颈（三指手背F1仅0.692）非注意力机制可解决。
  \end{minipage}
\end{table}
